\subsection{用关系来解释平坦性}

在本目中，A 表示一个环，E 表示一个右 A-模，F 表示一个左 A-模。

\(\mathbf{E} \otimes_{\mathbf{A}} \mathbf{F}\) 的每个元素都至少可以用一种方式写成 \(z = \sum_{i=1}^{n} e_i \otimes f_i\) 的形式，其中 \(e_i \in \mathbf{E}\) ， \(f_i \in \mathbf{F}\) 。下面的引理给出使这个和为零的一个条件：

引理 10. 设 \((f_{\lambda})_{\lambda \in \mathbf{L}}\) 是 \(\mathbf{F}\) 的一个生成元族， \((e_{\lambda})_{\lambda \in \mathbf{L}}\) 是 \(\mathbf{E}\) 的一个支集有限的元素族。为使 \(\sum_{\lambda \in \mathbf{L}} e_{\lambda} \otimes f_{\lambda} = 0\) ，必须且只需存在有限集 \(\mathbf{J}\) 、元素族 \((x_{j})_{j \in \mathbf{J}}\) （其元素属于 \(\mathbf{E}\) ）以及元素族 \((a_{j\lambda})\) ( \(j \in \mathbf{J}, \lambda \in \mathbf{L}\) )（其元素属于 \(\mathbf{A}\) ），它们具有下列性质：

\begin{enumerate}
\def\labelenumi{(\arabic{enumi})}
\item
  族 \((a_{j\lambda})\) 的支集有限；
\item
  \(\sum_{\lambda \in L} a_{j\lambda} f_{\lambda} = 0\) 对所有 \(j \in J\) 成立；
\item
  \(e_{\lambda} = \sum_{j\in J}x_ja_{j\lambda}\) 对所有 \(\lambda \in \mathbf{L}\) 成立。
\end{enumerate}

不严格地说，诸 \(e_{\lambda}\) 所成的组，必须是 A 中元素的某些组（它们是“诸 \(f_{\lambda}\) 之间的关系”）的以 E 为系数的线性组合。

考虑自由 A-模 \(\mathrm{A}_{s}^{(\mathrm{L})}\) 、它的典范基 \((u_{\lambda})\) 以及同态 \(g: \mathrm{A}_{s}^{(\mathrm{L})} \to \mathrm{F}\) ，它满足 \(g(u_{\lambda}) = f_{\lambda}\) （对所有 \(\lambda \in L\) ）；用 R 表示 g 的核，则（由于诸 \(f_{\lambda}\) 生成 F）有正合列

\[
\mathrm{R} \xrightarrow {i} \mathrm{A} _ {s} ^ {(L)} \xrightarrow {g} \mathrm{F} \longrightarrow 0
\]

其中 i 表示典范单射。由第 1 目的引理 1，我们导出正合列

\[
\mathrm{E} \otimes_ {\mathrm{A}} \mathrm{R} \xrightarrow {1 \otimes i} \mathrm{E} \otimes_ {\mathrm{A}} \mathrm{A} _ {s} ^ {(L)} \xrightarrow {1 \otimes g} \mathrm{E} \otimes_ {\mathrm{A}} \mathrm{F} \longrightarrow 0.\tag{13}
\]

现在， \(\mathbf{E} \otimes_{\mathbf{A}} \mathbf{A}_{s}^{(L)}\) 典范地等同于 \(\mathbf{E}^{(L)}\) ，族 \(e = (e_{\lambda}) \in \mathbf{E}^{(L)}\) 等同于 \(\sum_{\lambda \in L} e_{\lambda} \otimes u_{\lambda}\) （《代数学》第二章 §3 第 7 目命题 7 的推论 1）。为使这样的族属于 \(1_{\mathbf{E}} \otimes g\) 的核，必须且只需 \(\sum_{\lambda \in L} e_{\lambda} \otimes f_{\lambda} = 0\) 在 \(\mathbf{E} \otimes_{\mathbf{A}} \mathbf{F}\) 中成立；考虑到正合列 (13)，这等价于说 \(e\) 属于 \(1_{\mathbf{E}} \otimes i\) 的像，换言之，存在形如

\[
\sum_ {\lambda \in L} e _ {\lambda} \otimes u _ {\lambda} = \sum_ {j \in J} x _ {j} \otimes i (r _ {j})\tag{14}
\]

其中 \(x_{j} \in \mathbf{E}\) ， \(r_{j} \in \mathbf{R}\) ，且 \(\mathbf{J}\) 为有限集。写 \(i(r_{j}) = \sum_{\lambda \in \mathbf{L}} a_{j\lambda} u_{\lambda}\) ，则假设 \(r_{j} \in \mathbf{R}\) 蕴含关系 \(\sum_{\lambda \in \mathbf{L}} a_{j\lambda} f_{\lambda} = 0\) 对每个 \(j \in \mathbf{J}\) 成立；另一方面，关系式 (14) 蕴含 \(e_{\lambda} = \sum_{j \in J} x_{j} a_{j\lambda}\) 对每个 \(\lambda \in \mathbf{L}\) 成立（《代数学》第二章 §3 第 7 目命题 7 的推论 1），这就完成了证明。

命题 13. 为使 E 为 F-平坦的（第 2 目，定义 1），必须且只需满足下列条件：

\begin{enumerate}
\def\labelenumi{(\Alph{enumi})}
\setcounter{enumi}{17}
\tightlist
\item
  若 \((e_i)_{i \in I}\) 与 \((f_i)_{i \in I}\) 是 E 与 F 的元素的两个有限族，使得 \(\sum_{i \in I} e_i \otimes f_i = 0\) 在 E \(\otimes_{\mathbf{A}} \mathbf{F}\) 中成立，则存在有限集 J、元素族 \((x_j)_{j \in J}\) （其元素属于 E）与元素族 \((a_{jl})\) ( \(j \in J\) ， \(i \in I\) )（其元素属于 A），它们具有下列性质：
\end{enumerate}

\begin{enumerate}
\def\labelenumi{(\arabic{enumi})}
\item
  \(\sum_{i\in I}a_{ji}f_i = 0\) 对所有 \(j\in J\) 成立
\item
  \(e_i = \sum_{j \in J} x_j a_{ji}\) 对所有 \(i \in I\) 成立。
\end{enumerate}

假设 E 是 F-平坦的。设 \((e_{i})\) 与 \((f_{i})\) 是使 \(\sum_{i\in I}e_{i}\otimes f_{i}=0\) 在 \(E\otimes_{A}F\) 中成立的元素的有限族，并设 \(F'\) 为由诸 \(f_{i}\) 生成的 F 的子模。由于典范映射 \(E\otimes_{A}F'\to E\otimes_{A}F\) 是单射， \(\sum_{i\in I}e_{i}\otimes f_{i}=0\) 在 \(E\otimes_{A}F'\) 中同样成立，于是可对 E 与 \(F'\) 应用引理 10 ；这样便得到满足 (R) 中诸条件的族 \((x_{j})\) 与 \((a_{ji})\) 。

反之，假设条件 (R) 成立。设 \(\mathbf{F}'\) 为 \(\mathbf{F}\) 的一个子模，并设 \(y = \sum_{i\in I}e_i\otimes f_i\) 为典范映射 \(\mathbf{E}\otimes_{\mathbf{A}}\mathbf{F}'\to \mathbf{E}\otimes_{\mathbf{A}}\mathbf{F}\) 的核中一个元素。由于 (R) 成立，存在满足条件 (1) 与 (2) 的族 \((x_j)\) 与 \((a_{jl})\) 。我们断言，在 \(\mathbf{E}\otimes_{\mathbf{A}}\mathbf{F}'\) 中有

\[
y = \sum_ {i, j} x _ {j} a _ {j i} \otimes f _ {i} = \sum_ {j \in J} (x _ {j} \otimes \sum_ {i \in I} a _ {j i} f _ {i}) = 0.
\]

因而 \(\mathbf{E} \otimes_{\mathbf{A}} \mathbf{F}' \to \mathbf{E} \otimes_{\mathbf{A}} \mathbf{F}\) 是单射。

推论 1. 为使右 A-模 E 是平坦的，必须且只需下列条件成立：

(RP) 若 \((e_i)_{i\in I}\) 与 \((b_i)_{i\in I}\) 是 E 与 A 的元素的两个有限族，使得 \(\sum_{i\in I}e_i b_i = 0\) ，则存在有限集 J、元素族 \((x_j)_{j\in J}\) （其元素属于 E）与元素族 \((a_{ji})\) ( \(j\in J\) ， \(i\in I\) )（其元素属于 A），使得 \(\sum_{i\in I}a_{ji}b_i = 0\) 对所有 \(j\in J\) 成立，且 \(e_i = \sum_{j\in J}x_j a_{ji}\) 对所有 \(i\in I\) 成立。

条件 (RP) 就是命题 13 的条件 (R) 应用于模 \(\mathbf{F} = \mathbf{A}_s\) 的情形。

不严格地说，(RP) 表示：诸 \(b_{i}\) 之间每个以 E 为系数的“关系”，都是诸 \(b_{i}\) 之间以 A 为系数的“关系”的（以 E 为系数的）线性组合。

我们特别考虑 A 到环 B 的一个同态，它使 B 成为右 A-模。我们知道（第 3 目，命题 1），这等价于说该 A-模是平坦的，或等价于说它对每个左 A-模 \(A_{s}^{m}\) ( \(m \geqslant 1\) ) 是平坦的。取 E = B 与 \(F = A_{s}^{m}\) 应用命题 13 的条件 (R)，便得到下列条件：

推论 2. 为使环 B 是平坦的右 A-模，必须且只需它满足下列条件：

(RP') 每个由 B 的元素组成的解 \((y_{k})_{1\leqslant k\leqslant n}\) ，若是齐次线性方程组

\[
\sum_ {k = 1} ^ {n} y _ {k} c _ {k l} = 0 \quad (1 \leqslant i \leqslant m)\tag{15}
\]

（其系数 \(c_{kl}\) 属于 A）的解，都是线性组合

\[
y _ {k} = \sum_ {j = 1} ^ {q} b _ {j} z _ {j k} \quad (1 \leqslant k \leqslant n)\tag{16}
\]

其中系数为 \(b_{j} \in \mathbf{B}\) ，组合的对象是方程组 (15) 的解 \((z_{jk})_{1 \leqslant k \leqslant n}\) ，这些解由元素 \(z_{jk}\) （属于 \(\mathbf{A}\) ）组成。

